\FloatBarrier
\section{IR and Artifact Materials}
\label{sec:artifact}
The core value and index-expression forms are listed in
Tables~\ref{tab:value-expressions} and~\ref{tab:index-expressions}. Their
coordinate interpretation uses the mapping invariant in
Equation~\eqref{extended:eq:mapping-invariant}, complete disjoint partitions from
Section~\ref{sec:symbolic-axis-partitioning}, and ordered output interfaces. Normalized
terms have the denotation in Equation~\eqref{extended:eq:normalized-term}. Arithmetic
normalization is conditional on the selected laws and on the reduction-domain
and read-independence checks described in Section~\ref{sec:normalization}.
Algorithm~\ref{alg:mapping} gives the mapping construction, and
Appendix~\ref{app:mappings-to-attribute-programs} describes attribute construction and program
generation; these procedures do not establish completeness beyond the
mathematical fragment.

The retained evaluation directories contain exact benchmark sources,
original guards, generated programs, raw outcomes, timings, source hashes,
and implementation snapshots. \code{evaluation/current/README.md} documents
the earlier corpus, footprint, and original cost runs;
\code{evaluation/unified-construction/README.md} documents the current paired
full-inventory rerun and preserves the earlier baseline;
\code{evaluation/core-emission/README.md} documents emission-only results;
\code{evaluation/accepted-core/README.md} documents the accepted core path,
independent integer references, corruption/guard controls, and public
compile/load/application commands. Keeping these snapshots separate avoids
attributing the new path's successes to the historical corpus compiler.
These materials support reproduction of the reported experiments; they
are not a machine-checked formalization or a complete validation of every
supported operation and generated backend.
